% SI Text 2: the model in full; the long version's Section 3 with Table S1 (strategies).
\section{The model in full}\label{sec:model}

A funder faces a population of $n$ researchers over $T$ funding rounds. Researcher $i$ has capability $K_i$ and resources $R_i$. Initial capabilities and resources are drawn from Pareto distributions with tail parameters $\alpha_K$ and $\alpha_R$; smaller tail parameters produce more unequal fields, and a correlation parameter $\rho$ lets capability and resources be correlated across researchers. The funder knows these distributions; it never observes any researcher's capability or resources directly.

In each round, researcher $i$ produces research output\footnote{We are agnostic about the interpretation of output units: they may be interpreted as publications, as findings, or as units of scientific value.} at random, in discrete units, with expected output\footnote{Output in each round is drawn from a Poisson distribution with rate $\lambda_i$, independently across researchers and rounds given the state.}
\[
\lambda_i = A\,\Lambda(K_i, R_i),
\]
where $A$ is a productivity constant and $\Lambda$ is a production function that increases in both arguments but is limited by the scarcer of the two. Our main analysis takes $\Lambda$ to be proportional to the harmonic mean of the two inputs:
\[
\lambda_i = \frac{A K_i R_i}{K_i + R_i}.
\]
The harmonic mean makes capability and resources complements without making either an absolute requirement: a capable researcher with few resources still produces, by working more slowly or borrowing equipment, but no surplus of one input fully substitutes for a shortage of the other. \ref{app:technology} repeats our analysis for a family of production functions running from Cobb-Douglas, under which either input can substitute for the other, to Leontief, under which neither can.

Capabilities increase through research. Between rounds,
\[
K_i \leftarrow K_i + \epsilon\,\Lambda(K_i, R_i),
\]
where $\epsilon$ is the compounding rate: doing research builds expertise, and the same bottleneck that limits output limits growth. Growth requires both inputs, and capability never declines. Resources, by contrast, do not accumulate: resources in a round are baseline plus grant, $R_i = R_{i0} + g_i$. The baseline $R_{i0}$ recurs every round, representing the time and support a researcher has independent of the funder, while grants are consumed in the round they are given. Whatever lasts from a grant (skills, publications, standing) enters the model as capability.

The funder's total budget is fixed before any outcomes are observed. We parameterize it by a budget scale $b$: at $b = 1$, the budget equals the community's expected baseline resources for one round, and the budget scales linearly in $b$ (\ref{app:simulation}). The budget does not grow with the horizon: a longer horizon varies the timing of spending, never its amount. Strategies that do not plan ahead spend an equal share of the budget each round. Forward-looking strategies may spend on any schedule: each round, they re-plan the entire remaining budget over the remaining horizon and execute the current round's allocation.

The funder learns as it goes. It observes each researcher's track record, the output they have produced to date, as it accumulates, and two further signals are available. A resource signal, observed once at the start, gives a noisy estimate of each researcher's baseline resources (noise $\tau_R$), as one might obtain from institutional information. A review signal, drawn afresh in each round that a strategy uses it, gives a noisy estimate of capability (noise $\tau_K$), as one might obtain from proposal evaluation and expert judgment. The review signal's informativeness is governed by $\tau_K$: the smaller the noise, the more informative the review. Bayesian strategies update beliefs about every researcher after each round and allocate each dollar to the researcher whose expected output, given the funder's current beliefs, rises most with it.

We compare the strategies in Table~\ref{tab:strategies}. The Bayesian strategies differ along two dimensions: information, with or without the review signal, and planning, myopic (optimizing the current round's expected output) or forward-looking (optimizing expected output over the whole remaining horizon). Three baselines complete the set: no funding, against which every result is normalized; uniform funding, which splits each round's budget equally; and track-record funding, which allocates in proportion to the previous round's output. Seed grants and lotteries, which modify these strategies, are introduced in the main text.

\begin{table}[t]
\centering
\footnotesize
\setlength{\tabcolsep}{4pt}
\begin{tabular*}{\textwidth}{@{\extracolsep{\fill}}lll@{}}
\toprule
Strategy & Information & Planning \\
\midrule
No funding & none & none \\
Uniform & none & none \\
Track record & track record & none \\
Myopic, without review & track record + resource signal & current round \\
Forward, without review & track record + resource signal & remaining horizon \\
Myopic, with review & track record + resource signal + review & current round \\
Forward, with review & track record + resource signal + review & remaining horizon \\
\bottomrule
\end{tabular*}
\caption{The funding strategies compared.}
\label{tab:strategies}
\end{table}

Each of our conclusions rests on a proof or on simulations across the following parameter ranges; default values and the settings behind each figure are in \ref{app:simulation}. The tail parameters $\alpha_K$ and $\alpha_R$ range from 1.3 to 3.5. Empirical distributions of research productivity have Pareto tail parameters near 1 to 2 \citep{Lotka1926, Nicholls1989}; % [check Nicholls range]
our range covers these values and extends to much more equal fields. The budget scale $b$ ranges from 0.2 to 2, from a funder whose budget is small beside the community's own resources to one able to relieve most researchers' bottlenecks; the main text and the main text also examine larger budgets. The compounding rate $\epsilon$ ranges from near zero, where funding raises only current output, to 0.85, where this round's funding substantially raises later capability. The review noise $\tau_K$ ranges from 0.05, nearly perfect review, to 20, nearly uninformative review, a range wide enough to contain the noise implied by empirical estimates of review's predictive accuracy (main text). The correlation $\rho$ ranges from $-0.5$ to 0.8, and the horizon $T$ from one round to ten.

Randomness enters through the initial population draws, the signals, and realized output; capability growth follows expected output, so a researcher's accumulation depends on their capability and resources rather than on luck in a given round. We score each strategy by expected output and compare strategies on identical simulated populations, and we report each strategy's improvement over no funding in percent (\ref{app:simulation}).

With the model in place, the funder's problem is well defined: choose grants, given beliefs, to maximize expected total output over the horizon. Next we solve the case in which the funder has complete information regarding researchers' capabilities and resources.
